\section{Case Study}
\label{sec:case-study}

This section presents a case study of MOSS's evolution loop, using a controlled batch of \texttt{claweval}~\citep{ye2026claw} benchmark tasks as input to demonstrate empirically that the loop produces an effective harness-level fix. We feed four compliance-audit tasks to the loop as the failure batch, run the pipeline end-to-end (pre-loop baseline $\rightarrow$ stage pipeline $\rightarrow$ trial verification $\rightarrow$ host-daemon swap), and measure the iteration-1 candidate image against the baseline image on the same four tasks. The score difference, together with the qualitative change in agent output, is the evidence that the loop closes the gap from a controlled failure signal to a deployed agent improvement.

\subsection{Setup: Tasks and Baseline}
\label{subsec:case-setup}

The batch consists of four \texttt{claweval} tasks in the operations / compliance-audit category. T141zh / T142 (SLA compliance audit, Chinese / English) asks the agent to list P1 violation tickets and compute per-ticket SLA compliance against tiered rules served by a config mock service. T137zh / T138 (restock-chain check, Chinese / English) asks the agent to trace a multi-service chain---scheduler jobs, integration configs, inventory levels---and report which restock jobs are failing, which integration is broken, and which inventory is short. Each pair shares the same task across two language variants. We use these four tasks both as the input batch fed into evolution and as the test set on which the post-evolution candidate is re-scored---a controlled setup that lets the same grader witness both ends of the loop.

The agent under evolution is OpenClaw with DeepSeek V3.2~\citep{liu2025deepseek} as its underlying model. Baseline runs score 0.21--0.33 with a mean near 0.25 on the \texttt{claweval} grader's $[0,1]$ scale, well below its 0.75 pass threshold. Baseline trial transcripts show recurring failure modes across both task families: on the SLA-compliance pair, the agent typically lists only some of the relevant tickets, marks the others as ``response data incomplete, compliance status indeterminate'', and occasionally misattributes customer names between adjacent tickets; on the restock-chain pair, the response is similarly partial, with broken or missing chains between scheduler jobs, integrations, and inventories. These are the failure modes the evolution loop is asked to fix.

The \texttt{claweval} grader serves only as an external numeric witness for the reader; it is separate from MOSS's internal Task-Evaluate stage, which emits the qualitative keypoint matrix that drives the verdict (\S\ref{subsec:stage-pipeline}). Using a benchmark batch here trades realism for reproducibility---running the same four tasks pre- and post-evolution provides a controlled measurement that an organic user session could not anchor.

\subsection{What MOSS Does With the Batch}
\label{subsec:case-diag}

Given this input batch, MOSS runs the pre-loop baseline: the Task-Evaluate stage scores pre-captured baseline transcripts for each of the four tasks into a baseline keypoint matrix that locks the keypoint set for the rest of the loop. The matrix comes back weak or missing on tool sequencing, information extraction, and result reporting.

The Locate stage ingests the baseline transcripts and identifies a coverage gap in the harness's tool-result handling for multi-tool execution patterns. The agent under evolution routinely chooses one of the harness's generic execution paths over the semantic tools the mediator was designed around, and the mediator has no annotation branch for that path. A secondary parsing issue in the harness's dispatch-synthesis pipeline compounds the gap when the agent batches several lookups into a single shell construct, leaving downstream consumers with merged, partially-attributed outputs. Together these two surfaces explain the half-answers visible in the baseline trial transcripts.

The Plan stage translates the diagnosis into a two-surface fix inside the harness: an additional annotation branch that injects an explicit usage hint when the execution path returns a multi-entity payload, and a pre-call deny gate that blocks the specific batched-shell pattern so the agent is steered toward separate, individually-parsable calls. The Plan-Review stage approves the design without revisions. The Implement stage lands it as a single commit on the inner OpenClaw repository, modifying three files with 177 insertions and 1 deletion: a new annotation branch and supporting helpers in the tool-result mediator, an added pre-call check in the before-tool-call hook chain, and a new mediator test file exercising both surfaces. The change touches the harness itself rather than any text-mutable artifact, locating the fix where \S\ref{sec:intro} argues prior systems cannot reach.

The Code-Review stage approves, the host-daemon builds the candidate image, and trial workers replay the batch against it. Figure~\ref{fig:iter1-timeline} sketches the resulting iteration-1 timeline. The Task-Evaluate stage rescores the transcripts; the Verdict stage compares the new keypoint matrix against the baseline, observes broad lift across tool-sequencing and result-reporting keypoints, and the candidate converges.

\begin{figure}[ht]
\centering
\begin{tikzpicture}[
  font=\footnotesize,
  >={Stealth[length=2mm]},
  pipestep/.style={
    draw, thick, rounded corners=2pt, fill=blue!6,
    minimum width=2.0cm, minimum height=0.8cm, align=center
  }
]

% Row 1: reasoning stages 1, 2, 4 + Build
\node[pipestep] (s1) {\textcircled{\scriptsize 1}\,Locate};
\node[pipestep, right=4mm of s1] (s2) {\textcircled{\scriptsize 2}\,Plan};
\node[pipestep, right=4mm of s2] (s4) {\textcircled{\scriptsize 4}\,Implement};
\node[pipestep, fill=orange!12, right=4mm of s4] (build) {Build};

% Row 2: Trial + reasoning stages 6, 7 (left-aligned under Row 1)
\node[pipestep, fill=purple!10, below=9mm of s1] (trial) {Trial};
\node[pipestep, right=4mm of trial] (s6) {\textcircled{\scriptsize 6}\,Task-Evaluate};
\node[pipestep, right=4mm of s6] (s7) {\textcircled{\scriptsize 7}\,Verdict};

% Row 1 arrows
\draw[->, thick] (s1) -- (s2);
\draw[->, thick] (s2) -- (s4);
\draw[->, thick] (s4) -- (build);

% Wrap arrow: Build (end of row 1) → Trial (start of row 2)
\draw[->, thick] (build.south) -- ++(0,-3mm) -| (trial.north);

% Row 2 arrows
\draw[->, thick] (trial) -- (s6);
\draw[->, thick] (s6) -- (s7);

\end{tikzpicture}
\caption{Iteration-1 trace covering all stages MOSS executed; stage 3 (Plan-Review) and stage 5 (Code-Review) gates are elided for compactness.}
\label{fig:iter1-timeline}
\end{figure}

\subsection{Results: Iteration-1 Outcome}
\label{subsec:case-results}

On convergence, the host-daemon performs the in-place container swap described in \S\ref{subsec:swap}; for this benchmark run, the user-consent gate is auto-acknowledged so the evaluation pipeline can complete without interactive review. The converged candidate image is then re-run against the same four tasks, with grader scores reported in Table~\ref{tab:case-results}.

\begin{table}[ht]
\caption{Per-task \texttt{claweval} grader scores (mean of 3 trials per task) before and after the iteration~1 swap.}
\label{tab:case-results}
\centering
\begin{tabular}{lccc}
\toprule
Task & Baseline & Iter 1 & $\Delta$ \\
\midrule
T141zh\_sla\_compliance\_audit  & 0.3273          & 0.5330          & $+0.2057$ \\
T142\_sla\_compliance\_audit    & 0.2527          & 0.5453          & $+0.2926$ \\
T137zh\_restock\_chain\_check   & 0.2213          & 0.4567          & $+0.2354$ \\
T138\_restock\_chain\_check     & 0.2090          & 0.9049          & $+0.6959$ \\
\midrule
\textbf{mean}                    & \textbf{0.2526} & \textbf{0.6100} & $\boldsymbol{+0.3574}$ \\
\bottomrule
\end{tabular}
\end{table}

The four-task mean rises from 0.2526 to 0.6100. The largest gain is on T138, which climbs to 0.9049 with all three trials clearing the 0.75 pass threshold. The SLA pair (T141zh / T142) moves from the 0.25--0.33 range to roughly 0.53--0.55: the annotation path closes the largest defect, while per-ticket time-difference arithmetic and SLA-tier classification carry their own difficulty that one iteration does not resolve. T137zh shows the smallest gain at $+0.2354$, with one trial reaching pass and two below threshold.

The scores correspond to the agent's actual outputs. The SLA-compliance task transcripts that previously came back as the half-answers described in \S\ref{subsec:case-setup} now return per-ticket SLA-tier classifications and an aggregate summary (``2 of 6 P1 tickets violated SLA, affecting Customer B and Customer D, mean overrun 2.3h''); the restock-chain transcripts likewise return complete chains. The model and the task definition are unchanged; the harness change is the proximate cause.
